% ---
% titre: Réseaux sociaux des jeunes en Suisse (construction d'un diagramme en colonne)
% theme: FA
% chapitre: Diagrammes
% sous_chapitre: Réaliser des diagrammes
% niveau: [10e]
% type: EVAL
% tags: [c-diagramme-colonnes,c-proportionnalite,c-echelle,o-regle,p-question-TS]
% difficulte: 1
% corrige: true
% source: JAMESfocus 2025 (ZHAW), données JAMES 2024
% ---

\question[6]
L'étude suisse JAMES (2024) a interrogé plus de 1\,000 jeunes de 12 à 19 ans sur leur fréquence d'utilisation des réseaux sociaux. Le tableau ci-dessous indique la part d'entre eux qui se rend sur chaque plateforme plusieurs fois par semaine ou plus.

\vspace{40pt}\def\arraystretch{3}
\begin{center}Part des jeunes utilisant chaque plateforme\end{center}
\begin{tabularx}{\linewidth}{|*{7}{>{\centering\arraybackslash}X|}}
\hline
\rowcolor{gray!20}\textbf{Instagram} & \textbf{TikTok} & \textbf{BeReal} & \textbf{Pinterest} & \textbf{X/Twitter} & \textbf{Reddit} & \textbf{Facebook}\\
\hline
83\,\% & 69\,\% & 32\,\% & 26\,\% & 8\,\% & 5\,\% & 4\,\%\\
\hline
\end{tabularx}
\def\arraystretch{1}

\vspace{40pt}\textbf{Représente précisément ces données par un diagramme en colonnes.}


\vfill
Titre du diagramme :\quad \sol{\rule{0.75\linewidth}{0.4pt}}{Réseaux sociaux utilisés plusieurs fois par semaine par les jeunes en Suisse \textbf{(0.5pt)}}

\vfill{\footnotesize\raggedleft Espace pour ton diagramme (au crayon et à la règle)\par}

\begin{center}
\begin{tikzpicture}
  % correction dans le calque d'arrière-plan : les axes restent toujours noirs par-dessus
  \begin{pgfonlayer}{background}
    \begin{scope}[color=\colSolution]
      \foreach \i/\nom/\v in {0/Instagram/83, 1/TikTok/69, 2/BeReal/32, 3/Pinterest/26, 4/X\slash Twitter/8, 5/Reddit/5, 6/Facebook/4}{
        % colonne (largeur 1 cm, espacement 1 cm : alignée sur le quadrillage)
        \fill[\colSolution!15] (2.5+2*\i,1) rectangle (3.5+2*\i,1+\v/10);
        \draw[thick] (2.5+2*\i,1) rectangle (3.5+2*\i,1+\v/10);
        % zone de tolérance : hauteur exacte à plus ou moins 1 mm (1 cm = 10 %)
        \fill[\colSolution!50] (2.5+2*\i,1+\v/10-0.1) rectangle (3.5+2*\i,1+\v/10+0.1);
        \node[font=\scriptsize] at (3+2*\i,0.6) {\nom};
      }
    \end{scope}
  \end{pgfonlayer}
  % quadrillage (calque principal, donc toujours visible, même sous la correction)
  % carreaux de 0,5 cm, trait un peu plus marqué tous les 1 cm (= 10 %)
  \draw[step=0.5cm, gray!35, very thin] (2,1) grid (16,10);
  \draw[step=1cm, gray!60, very thin] (2,1) grid (16,10);
  % axes (origine en (2,1) ; 1 cm = 10 %)
  \draw[thick, ->] (2,1) -- (2,10.6);
  \draw[thick, ->] (2,1) -- (16,1);
  \node[anchor=west, font=\small] at (2.2,10.5) {Part des jeunes (\%)};
  \node[font=\small] at (15,0.0) {Plateformes};
  % graduations : traits tous les 0,5 cm, valeurs tous les 1 cm
  \foreach \p in {0,10,...,90}{
    \draw[thick] (1.85,1+\p/10) -- (2.15,1+\p/10);
    \node[left, font=\small] at (1.8,1+\p/10) {\p};
  }
  \foreach \p in {5,15,...,85}{
    \draw (1.92,1+\p/10) -- (2.08,1+\p/10);
  }
\end{tikzpicture}
\end{center}

\newpage\,\vfill{\footnotesize\raggedleft Espace pour ta démarche et tes calculs\par}
\begin{solutionorgrid}[610pt]

\textbf{1\,pt} pour le nom des plateformes sous les colonnes (ou dans les colonnes)\\
\textbf{0,5\,pt} pour les colonnes de même largeur, régulièrement espacées, au crayon et à la règle

\vspace{10pt}
\textbf{4\,pts} si les extrémités supérieures des sept colonnes sont dans les zones grisées\\
\textbf{3\,pts} si les extrémités supérieures de cinq ou six colonnes sont dans les zones grisées\\
\textbf{2\,pts} si les extrémité supérieure de trois ou quatre colonnes sont dans les zones grisées\\
\textbf{1\,pt} si les extrémités supérieures d'une ou deux colonnes sont dans les zones grisées

\vspace{10pt}
Déduire 1 pt en présence d'un diagramme cartésien (le tracé des colonnes est manquant).
\end{solutionorgrid}